\documentclass[aspectratio=169,10pt]{beamer}

\usetheme{default}
\usefonttheme{professionalfonts}
\setbeamertemplate{navigation symbols}{}
\setbeamertemplate{footline}{\hfill\scriptsize\insertframenumber/\inserttotalframenumber\hspace{2mm}\vspace{1mm}}

\usepackage{amsmath,amssymb,bm}
\usepackage{mathtools}
\usepackage{xcolor}
\usepackage[most]{tcolorbox}

\definecolor{deepblue}{RGB}{0,70,150}
\definecolor{softblue}{RGB}{232,242,255}
\definecolor{softgreen}{RGB}{235,250,238}
\definecolor{softorange}{RGB}{255,245,230}
\definecolor{darkred}{RGB}{160,35,25}
\definecolor{darkgreen}{RGB}{20,120,60}

\setbeamerfont{frametitle}{series=\bfseries,size=\Large}
\setbeamercolor{frametitle}{fg=black}
\setbeamertemplate{itemize item}{\color{deepblue}$\blacktriangleright$}
\setlength{\abovedisplayskip}{3pt}
\setlength{\belowdisplayskip}{3pt}

\newtcolorbox{bluebox}[1][]{colback=softblue,colframe=deepblue,boxrule=0.8pt,arc=1.5mm,left=1.5mm,right=1.5mm,top=0.8mm,bottom=0.8mm,#1}
\newtcolorbox{greenbox}[1][]{colback=softgreen,colframe=darkgreen,boxrule=0.8pt,arc=1.5mm,left=1.5mm,right=1.5mm,top=0.8mm,bottom=0.8mm,#1}
\newtcolorbox{orangebox}[1][]{colback=softorange,colframe=darkred,boxrule=0.8pt,arc=1.5mm,left=1.5mm,right=1.5mm,top=0.8mm,bottom=0.8mm,#1}

\begin{document}

\begin{frame}{Module 1: Convert particle transport into source maps}
\small
\vspace{-2mm}
\begin{bluebox}
\centering\small
\textbf{Goal:} convert microscopic Geant4 transport events into continuum source terms for the later 2D multiphysics modules.
\end{bluebox}

\vspace{1mm}
\begin{columns}[T,totalwidth=\textwidth]
\begin{column}{0.50\textwidth}
\textbf{Continuum definition}
\begin{equation*}
\boxed{\dot{q}_{\mathrm{rad}}(\bm r,t)=\frac{\partial E_{\mathrm{dep}}}{\partial V\,\partial t}}
\end{equation*}
\small
\begin{itemize}\setlength\itemsep{0.5mm}
\item \(\dot{q}_{\mathrm{rad}}\): radiation-induced volumetric energy-deposition rate.
\item Units: \(\mathrm{W/m^3}=\mathrm{J/(m^3s)}\).
\item It is an \emph{initial forcing map}, not automatically all heat.
\end{itemize}
\end{column}

\begin{column}{0.48\textwidth}
\textbf{Scored mesh / histogram form}
\begin{equation*}
\boxed{\dot{q}_{\mathrm{rad},i}\approx
\frac{1}{\Delta V_i\Delta t}\sum_{k\in i}E_{\mathrm{dep},k}}
\end{equation*}
\small
\begin{itemize}\setlength\itemsep{0.5mm}
\item \(i\): mesh cell or scoring bin.
\item \(E_{\mathrm{dep},k}\): energy deposited by step/event \(k\).
\item \(\Delta V_i\), \(\Delta t\): cell volume and scoring time window.
\end{itemize}
\end{column}
\end{columns}

\vspace{1mm}
\begin{orangebox}
\centering\small
Downstream partition:
\[
\dot{q}_{\mathrm{rad}}(\bm r,t)\rightarrow
\left\{\dot{q}_{\mathrm{th}}(\bm r,t),\;G_{eh}(\bm r,t),\;G_j(\bm r,t)\right\},
\]
where \(\dot{q}_{\mathrm{th}}\) is thermalized heating, \(G_{eh}\) is electron-hole generation, and \(G_j\) is defect generation.
\end{orangebox}
\end{frame}

\begin{frame}{Module 1: Map 3D Geant4 scoring to the 2D project model}
\small
\vspace{-2mm}
\begin{bluebox}
\centering\small
The project keeps in-plane device coordinates \((x,y)\) and collapses the out-of-plane/depth direction \(z\) by averaging.
\end{bluebox}

\vspace{1mm}
\begin{columns}[T,totalwidth=\textwidth]
\begin{column}{0.52\textwidth}
\textbf{Depth-averaged 2D source}
\begin{equation*}
\boxed{
\dot{q}_{\mathrm{rad}}^{xy}(x,y,t)
=\frac{1}{L_z}\int_{z_{\min}}^{z_{\max}}
\dot{q}_{\mathrm{rad}}(x,y,z,t)\,dz}
\end{equation*}
\small
\begin{itemize}\setlength\itemsep{0.5mm}
\item \(L_z=z_{\max}-z_{\min}\): collapsed depth.
\item Units remain \(\mathrm{W/m^3}\), so the result can enter 2D PDEs as a volumetric source.
\end{itemize}
\end{column}

\begin{column}{0.46\textwidth}
\textbf{Binned 3D histogram reduction}
\begin{equation*}
\boxed{
\dot{q}_{\mathrm{rad},ij}^{xy}(t)
=\frac{1}{L_z}\sum_k\dot{q}_{\mathrm{rad},ijk}(t)\Delta z_k}
\end{equation*}
\small
For uniform depth bins,
\begin{equation*}
\boxed{
\dot{q}_{\mathrm{rad},ij}^{xy}(t)
=\frac{1}{N_z}\sum_{k=1}^{N_z}\dot{q}_{\mathrm{rad},ijk}(t)}
\end{equation*}
\end{column}
\end{columns}

\vspace{1mm}
\begin{greenbox}
\small
\textbf{Preserves:} average cross-sectional forcing seen by Modules 2--6.\quad
\textbf{Loses:} out-of-plane localization, peak 3D hot spots, and event-to-event asymmetry in \(z\).
\end{greenbox}

\vspace{1mm}
\begin{orangebox}
\centering\small
Downstream forcing: \(\dot{q}_{\mathrm{rad}}^{xy}\rightarrow\{G_j,\dot{q}_{\mathrm{th}},G_{eh}\}\rightarrow\{C_j,T,\phi,n,p,\bm J_n,\bm J_p,y_{\mathrm{device}}\}\).
\end{orangebox}
\end{frame}

\end{document}
